Large language models are increasingly used as stand-ins for people in product evaluation, social simulation, and agent design. The central question is not whether these models can produce human-like text, but whether they can simulate the distribution of human outcomes once we condition them on latent factors such as demographics or persona. We propose a simple theory of simulation with three requirements: local fidelity to known human-style judgments, responsiveness to latent factors, and coverage of plausible behavioral variation. We test this framework with real API calls to \gptfourone on May 5, 2026 using two tasks: a 90-item \socialiqa social-judgment benchmark and a Turing-Experiment-style ultimatum game with 15 synthetic participants per condition. We compare unconditioned, demographic-conditioned, and persona-conditioned prompting while holding the base model fixed.

Persona conditioning improves \socialiqa accuracy from 0.744 to 0.789 and raises the number of distinct ultimatum policies from 2 to 5. However, all conditions still accept unfair offers at unrealistically high rates, with mean acceptance thresholds between 2.40 and 2.67 on a 1--9 offer scale. The added diversity therefore broadens behavior only within a narrow cooperative family. Our main conclusion is that current frontier LLMs are controllable but compressed behavioral simulators: they can support local fidelity and some latent-factor responsiveness, but they do not yet simulate a true superset of human outcomes. This distinction matters for any application that relies on subgroup structure or tail behavior.
